\section{A verified bounded arithmetic construction}
\label{app:arithmetic}

This appendix proves Lemma~\ref{lem:basic-arithmetic}. The arithmetic
identities follow the conjunction-only reductions of
\citet[Lemmas C--G and Figures 2--3]{AbrahamsenMiltzow2019}; we give the
details needed for rational equivalence and use compactness and continuity
for the range choices. In particular, we do not use its Boolean
normal-form step or quantitative range lemma.

\subsection{Unique polynomial evaluation and scaling}

Clear positive rational denominators in a basic closed description of
$S$. Evaluate its integer polynomials by a finite arithmetic circuit with
equations of the forms
\[
 z=1,\qquad x+y=z,\qquad xy=z,
\]
and require the appropriate output coordinates to be zero or nonnegative.
The constant zero is fixed by $0+1=1$, and a negative value is introduced
by $x+(-x)=0$. Integer constants are formed by repeated addition and
negation. A zero output can be imposed by $z+0=0$.
Each circuit coordinate is a uniquely determined polynomial in the
original coordinates. The solution set of this conjunctive system is
therefore the graph of a polynomial map on $S$, with coordinate projection
as inverse. It is compact. Write $M\ge1$ for a bound on the absolute
values of all circuit coordinates on that set; the empty case can be
represented directly by the inconsistent bounded equation $x+x=x$.

Choose a sufficiently small dyadic number $\delta=2^{-\ell}>0$, as
specified below, and a dyadic $\varepsilon=\delta 2^{-k}$ with
$k\ge0$ and $\varepsilon M\le\delta$. Replace each circuit variable
$x$ by $w_x=\varepsilon x$. Introduce a variable $d=\delta$, a zero
variable via $0+d=d$, and the halving chain from $d$ to $\varepsilon$.
These constants are all uniquely determined. Replace the circuit equations by
\begin{align*}
 x=1&:\quad w_x+0=\varepsilon,\\
 x+y=z&:\quad w_x+w_y=w_z,\\
 xy=z&:\quad w_xw_y=q,\quad \varepsilon w_z=q.
\end{align*}
Here $q$ is a new variable for that multiplication. Nonnegativity becomes
$w_x\ge0$. Every solution of this new system recovers a solution of the
old one by $x=w_x/\varepsilon$; the converse extension is unique and
has $q=\varepsilon^2z$. Thus all its variables lie in $[-\delta,\delta]$
when $\delta\le1$. This is a consequence of the equations and the
compact original set, not an additional bound that could remove solutions.
We need only existence of $\varepsilon$, and make no complexity claim
about the length of its halving chain.

\subsection{Shifted addition and multiplication}

All variables in the final construction have the common interval
$[1/2,2]$. A constant $c$ in the following displays denotes a dedicated
variable fixed to that value by the construction. First fix the variable $c_1$ by
$c_1c_1=1$, an inversion equation of $\ETRINV$; positivity excludes
the root $-1$. We then denote this fixed variable by $1$. The additions
\[
 \tfrac12+\tfrac12=1,\quad
 1+\tfrac12=\tfrac32,\quad
 \tfrac34+\tfrac34=\tfrac32,\quad 1+1=2
\]
fix the other displayed constants. Fix $2/3$ by the inversion
$(3/2)(2/3)=1$.

For each small variable $s$, retain $A_s=1+s$ and introduce
\begin{equation}
\label{eq:shift-coordinates}
 A_s+\tfrac12=C_s,\qquad B_s+\tfrac34=C_s.
\end{equation}
These uniquely enforce $C_s=s+3/2$ and $B_s=s+3/4$.
An addition $s+t=u$ is replaced by $B_s+B_t=C_u$.
For a product $st=u$, introduce the following equations, with fresh
$m,f,g,h$:
\begin{equation}
\label{eq:shift-product}
 \begin{aligned}
 A_sA_t&=m,& m+\tfrac12&=f,& g+B_t&=f,\\
 g+\tfrac34&=h,& B_u+B_s&=h.
 \end{aligned}
\end{equation}
Successive substitution gives
\[
 m=1+s+t+st,\quad f=\tfrac32+s+t+st,\quad
 g=\tfrac34+s+st,\quad h=\tfrac32+s+st.
\]
The last equation is precisely $u=st$.
To impose $s\ge0$, add $J_s+1/2=A_s$; its built-in bound
$J_s\ge1/2$ is exactly the desired inequality.

It remains to impose $d=\delta$. Construct $D_i=1-2^{-i}$ for
$i=1,\ldots,\ell$: start with $D_1=1/2$, and use
\[
 D_i+1=E_i,\qquad D_{i+1}+D_{i+1}=E_i.
\]
All these constants lie in $[1/2,2]$. With $\Delta=D_\ell=1-\delta$,
the equation $A_d+\Delta=2$ fixes $d=\delta$ exactly. The small
system's other equations, including its halving chain, are replaced by
the addition and product rules above. For $\delta$ sufficiently small,
all nonconstant coordinates in these rules are in $(1/2,2)$, except
$J_s=1/2+s$, whose range $[1/2,1/2+\delta]$ follows from $s\ge0$.
The only multiplications still present have their two operands near $1$.
Their replacement is described next.

\subsection{Products from squares, and squares from reciprocals}

In the next two displays, each assignment introduces a fresh variable,
except that the final output is identified with the desired existing
output coordinate. An assignment $v=u-w$ means the addition equation
$v+w=u$, and $v=u/2$ means $v+v=u$. Thus subtraction and halving
introduce no new allowed operation. The order of the assignments makes
all auxiliary values unique.

For inputs $a,b$ near $1$, the following construction uses only additions
and three squarings:
\begin{equation}
\label{eq:product-square-chain}
\begin{aligned}
 a'&=a+\tfrac12,& b'&=b+\tfrac12,&
 h_a&=a'/2,& h_b&=b'/2,\\
 t&=h_a+h_b,& r&=t-\tfrac12,& s&=r^2,& u&=s+\tfrac12,\\
 a_2&=a^2,& b_2&=b^2,& a_3&=a_2+\tfrac12,& b_3&=b_2+\tfrac12,\\
 a_4&=a_3/2,& b_4&=b_3/2,& c&=a_4+b_4,& d&=c/2,\\
 e&=u-d,& f&=e+e,& o&=f-\tfrac12.&&
\end{aligned}
\end{equation}
Indeed $r=(a+b)/2$, $d=(a^2+b^2+1)/4$, and
\[
 o=2\left(\frac{(a+b)^2}{4}+\frac12
             -\frac{a^2+b^2+1}{4}\right)-\frac12=ab.
\]
At $a=b=1$, every variable in this display is one of $3/4,1,3/2$.
In particular all lie strictly inside the allowed interval, and each
of the three squared inputs $r,a,b$ is near $1$.

For an input $a$ near $1$, replace its squaring by
\begin{equation}
\label{eq:square-inverse-chain}
\begin{aligned}
 b&=a+\tfrac12,& b'&=1/b,& a'&=1/a,\\
 c&=a'+\tfrac12,& d&=c-\tfrac23,& e&=d+\tfrac12,\\
 f&=e-b',& g&=f+f,& h&=g-\tfrac23,\\
 i&=1/h,& j&=i-\tfrac12,& k&=j+\tfrac34,\\
 l&=b/2,& o&=k-l.&&
\end{aligned}
\end{equation}
Each reciprocal assignment $v=1/u$ is the allowed equation $uv=1$.
The central identity is
\[
 h=2\left(\frac1a+\frac13-\frac1{a+1/2}\right)-\frac23
   =\frac1{a(a+1/2)}.
\]
Consequently $i=a^2+a/2$ and
$o=i-1/2+3/4-(a+1/2)/2=a^2$.
At $a=1$, the variables belong to
\[
 \{2/3,3/4,5/6,1,4/3,3/2,7/4\}\subset(1/2,2).
\]
All denominators are nonzero there.

For completeness, the range choices can be made before scaling the
original circuit. The finitely many rational functions in
\eqref{eq:square-inverse-chain} are continuous with nonzero denominators
at $a=1$, so they all remain in $(1/2,2)$ on one neighborhood of $1$.
The finitely many polynomial functions in
\eqref{eq:product-square-chain} are continuous at $(1,1)$, and its
three squared inputs approach $1$ there. Shrink the product-input
neighborhood so that all its variables stay in $(1/2,2)$ and its squared
inputs enter the reciprocal gate's neighborhood. Finally choose a dyadic
$\delta>0$ small enough for the shifted-product inputs $A_s,A_t$ and
all the shifted auxiliary coordinates above. This choice is uniform over
the finite gate types; it does not depend on the number of circuit gates.
The constants constructed by the midpoint and halving chains already
have their stated exact ranges and need no continuity argument.

Every remaining constraint is now an addition or inversion with variables
in $[1/2,2]$. On each original solution the formulas give one valid
extension. Conversely, every bounded solution forces the constants,
reciprocal gates, square gates, and shifted equations by the identities
above. It therefore recovers the small system, and then the original
circuit and point of $S$. Each added coordinate is a rational function
of the original point with nonzero denominators on $S$; recovery is
$t_i=(A_{w_{t_i}}-1)/\varepsilon$.
Both maps are continuous, and this proves
Lemma~\ref{lem:basic-arithmetic}, including its designated-coordinate claim.

\subsection{Why the input must be conjunctive}

The general rational-equivalence assertion in
\citet[Theorem 1]{AbrahamsenMiltzow2019} is broader than the part used
here. Its Lemma A replaces weak inequalities by equations with
nonnegative auxiliary variables inside their Boolean branches.
For example, applying that rule and eliminating a disjunction by a product
turns $x\ge0$ or $y\ge0$ into
\[
 (x-u)(y-v)=0,\qquad u\ge0,\quad v\ge0.
\]
At $(x,y)=(-1/2,1/2)$, every $u\ge0$ with $v=1/2$ satisfies these
conditions. Thus the auxiliaries are not uniquely determined; adding
the box $[-1,1]^2$ to the original variables does not repair that failure.
The displayed replacement preserves existence of a solution,
but it does not give the claimed rational bijection. The basic-closedness obstruction proved in
Section~\ref{sec:algebraic} shows that unrestricted rational equivalence
cannot be restored by choosing a different Boolean preprocessing rule.
Our polynomial evaluation starts from a conjunction and defines every
auxiliary globally, so it has no inactive branches.
